% ================= DISCUSSION =================
\section{Discussion}
\label{sec:discussion}

\paragraph{What the hypothesis got right and wrong.} The proposal expected the classical estimator
to be weakest on text because sparse strokes give a weak blur signature. The controlled sweep
separates two effects the hypothesis had merged: \emph{content} and \emph{region size}. At full
frame, natural images are easy; cut to a plate-sized window they become the hardest class, while
printed text in the same window is the easiest --- its high-contrast, broadband strokes are exactly
what produces deep, regular nulls. Plates fall in between because the window also holds bumper and
background. What drives the failures is the number of samples that carry the null pattern (region
size), the noise level, and short blurs whose first null lies far out in the spectrum. Plates are
hard mainly because they are small.

\paragraph{Why edge enhancement fails and band limiting works.} The proposal's pre-processing idea
was to strengthen the strokes so the nulls stand out. In the frequency domain this is a high-pass
filter: it raises exactly the band where noise, JPEG blocking and resampling artefacts live and
where the sinc nulls are already the weakest, and it adds structure of its own. The blur
information of a long blur sits at low frequencies (its nulls are spaced $1/L$ apart), so the useful
operation is the opposite one --- discard the upper band. The multi-band search keeps the benefit
for short blurs by still trying the full band and letting the directionality of the dip decide.
The one edge-type operator that helped on synthetic data, a mild Gaussian high-pass, is consistent
with this picture: it flattens the steep low-frequency fall-off of the spectrum without boosting
the noise band, whereas the Laplacian multiplies the spectrum by $|\omega|^2$.

\paragraph{The gate is what makes the pipeline usable.} On real images the dominant risk is not a
slightly wrong kernel but a wrong \emph{model}: defocus, curved or non-uniform shake and compression
cannot be explained by one straight line, yet the estimator always returns one. Dip depth alone
cannot tell these cases apart. The test that does is physical: a uniform motion blur must give the
same estimate in every part of the region. This is cheap (four more estimates), needs no training
and has an interpretable refusal reason.

\paragraph{Limitations.}
\begin{itemize}
  \item One spatially uniform straight-line kernel per image (as scoped in the proposal). The
        K\"ohler results show how quickly this breaks on curved camera shake --- even the best line
        fit loses 1\,dB --- and the gate does not detect that failure, because a curved trajectory
        is still consistent across image halves.
  \item LPBlur has no text labels; the reference is OCR of the paired sharp image, so plates whose
        sharp image is misread are excluded, and the reader's own errors add noise.
  \item The gate is calibrated on a mixture of synthetic and real development data. Its thresholds
        reflect that mixture; a deployment with a very different blur distribution would need
        recalibration (the calibration script does this from labelled examples).
  \item The OCR read-out is a single off-the-shelf recogniser; absolute accuracies depend on it, the
        comparisons between systems much less so.
\end{itemize}

\paragraph{Future work.} Natural extensions within the classical toolbox are a two-segment or
curved kernel model with the same spectral evidence, local (patch-wise) estimation to handle
non-uniform shake, and using the split-half estimates not only to refuse but to detect where the
blur changes.
